\section{Controlling triangles by splitting vertices}
\label{sec:triangles}

We prove Proposition~\ref{prop:triangle-bound} by repeatedly reducing the
neighborhood weights in the level graphs.  The main construction splits a
vertex into copies of the same weight, distributes its incident edges among
those copies, and preserves most of the triangle mass.  Two features must hold together: the
edges of a typical triangle must meet at the same copy of each corner, and
each copy must have a much smaller neighborhood weight.  Random walks inside
neighborhoods supply labels that achieve both properties.  The argument
adapts the triangle estimate of~\cite[Sections~3--5]{OpenAIIndependence},
while retaining the bounds on radius and participation in each part
needed for the weight adjustments here.

Throughout this section, a graph at level $b$ has a homomorphism into $H_b$
that is injective on directed edges.  Its vertices inherit the weights $q$,
prices, parts, and low, middle, or high classification of their images.
Corollary~\ref{lem:local-cross-transfer} therefore makes
\eqref{eq:local-cross} available on each such graph, with its stated ball and
part conditions.  Deleting vertices or edges preserves these properties.
The splitting construction will preserve them as well.

\subsection{Removing triangles of small total mass}
\label{subsec:triangle-preparation}

Let $F_b$ denote the current graph at level $b$, initially $H_b$, and put
\[
 T=\sum_b b\sum_{\{a,u,v\}\text{ a triangle of }F_b}q_aq_uq_v.
\]
Thus initially $T=T_0$.  A triangle of $F_b$ projects to a triangle of
$H_b$ with three distinct vertices and hence has a middle vertex.  We shall
use the following bound at every stage:
\begin{equation}
 \sum_b b\sum_{\substack{a\in V(F_b)\\a\text{ middle}}}
             q_a q\bigl(N_{F_b}(a)\bigr)\le 2M_x.
 \label{eq:middle-edge-charge}
\end{equation}
Indeed, for an edge $au$ the quantity $bq_aq_u$ is at most the contribution
of its image edge to $M_x$, since the coefficient $\kappa$ on that image is
at least one.  An original vertex is middle at exactly one level, and a
directed edge at that level has at most one preimage.  Each original edge
can consequently be charged at most once from each endpoint.

Suppose that all current neighborhood weights are at most $e^\ell$, where
$32\le\ell\le\ell_0$.  First delete, at every level, the vertices whose
inherited prices satisfy $x_i/b<\ell^{-4}$.  To bound the triangle mass
lost, fix a middle vertex $a$ before these deletions, let
$D=N_{F_b}(a)$, and let $S$ be the set of vertices of $D$ being deleted.
The set $S\cup D=D$ lies in a ball of radius one and has at most one vertex
in each part.  Also $q(S),q(D)\le e^\ell$, while
$\alpha_i,\beta_i\le\ell^{-4}$ on $S$.  Applying
\eqref{eq:local-cross} with $h=\ell$ gives
\[
 e_q(S,D)\le C_1\bigl(\ell^{-3}q(D)+e^{-3\ell}\bigr).
\]
In fact $e_q(S,D)=O_r(q(D)/\ell^2)$: when $q(D)>\ell^{-2}$ the
exponential term is absorbed by $q(D)/\ell^2$; when
$q(D)\le\ell^{-2}$, use instead $e_q(S,D)\le q(D)^2$.
A deleted triangle has a middle vertex, which is
not itself deleted, and an edge from $S$ to $D$ in that middle vertex's
neighborhood.  Multiplying by $bq_a$ and summing therefore bounds the loss
in $T$ by $O_r(M_x/\ell^2)$, using
\eqref{eq:middle-edge-charge}.

After these deletions, the total edge mass over the levels satisfies
\begin{equation}
 \sum_b b\sum_{uv\in E(F_b)}q_uq_v
                  \le O(\log\ell)M_x.
 \label{eq:dyadic-edge-mass}
\end{equation}
To see this, write $x_{\min}$ for the smaller of the two prices on an
original edge.  If that edge has a preimage at level $b$, then
\[
 b\ell^{-4}\le x_{\min}<2b.
\]
The first inequality follows from the deletion rule, and the second from
the absence of edges with two high endpoints.  There are only
$O(\log\ell)$ dyadic choices for $b$.  At each such level, the edge has
at most one preimage, whose mass $bq_uq_v$ is at most the original edge's
contribution to $M_x$.

Next delete an edge whenever the weight of its current common
neighborhood is less than $\ell^{-1}$, continuing until no such edge
remains.  Deleting $uv$ destroys triangle mass at most
$bq_uq_v/\ell$.  Thus \eqref{eq:dyadic-edge-mass} bounds this further
loss by $O((\log\ell)M_x/\ell)$.  Discard isolated vertices as well.

Fix a nonempty level graph $F=F_b$ after these operations.  All
neighborhoods in the next two subsections are in $F$.  Set
\begin{equation}
 \begin{aligned}
 S_u&=q(N(u))\in[\ell^{-1},e^\ell],\\
 c_{uv}&=q(N(u)\cap N(v))\ge\ell^{-1}
                       &&(uv\in E(F)),\\
 h_u&=\sum_{v\in N(u)}q_vc_{uv}\ge S_u/\ell.
 \end{aligned}
 \label{eq:prepared-level}
\end{equation}
The lower bound on $S_u$ follows because $u$ has a neighbor and every edge
has common-neighborhood weight at least $\ell^{-1}$.

\begin{lemma}[Splitting a level graph]
\label{lem:triangle-splitting}
Let $32\le\ell\le\ell_0$, and let $F$ be a nonempty finite graph with no isolated
vertices, with a homomorphism into some $H_b$ injective on directed edges.
Give its vertices the inherited weights, prices, and parts, and suppose
that \eqref{eq:prepared-level} holds.  There is a graph $F'$ obtained by
splitting vertices and retaining at most one copy of each edge of $F$,
with unchanged vertex weights, such that every neighborhood of $F'$ has
weight at most $e^{\sqrt\ell}$ and
\begin{equation}
 \sum_{\{a,u,v\}\text{ a triangle of }F'}q_aq_uq_v
 \ge \left(1-C_2\frac{\log\ell}{\sqrt\ell}\right)
       \sum_{\{a,u,v\}\text{ a triangle of }F}q_aq_uq_v,
 \label{eq:split-retention}
\end{equation}
where $C_2$ depends only on $r$.  The projection $F'\to F$ is a
homomorphism injective on directed edges.
\end{lemma}

We prove the lemma in the next two subsections.  At a vertex $u$, the
construction will attach a label in $N(u)$ to each incidence $(u,v)$;
edges with the same label use the same copy of $u$.  A triangle survives
only if its two labels agree at each corner.  We first construct
probability laws whose labels can usually agree, then impose a condition
on the reverse transition probabilities to bound the weight at each
copy.

\subsection{Entropy bounds for walks in neighborhoods}
\label{subsec:neighborhood-entropy}

Let $t$ be the unordered triangle mass of $F$.  It is positive by
\eqref{eq:prepared-level}.  Choose an ordered triangle $(u,v,z)$ with
probability $q_uq_vq_z/(6t)$.  The law of its first two vertices and the
conditional law of its second vertex are
\begin{equation}
 \mu(u,v)=\frac{q_uq_vc_{uv}}{6t},\qquad
 \pi_u(v)=\frac{q_vc_{uv}}{h_u}.
 \label{eq:triangle-laws}
\end{equation}
In particular, the marginal law of $u$ is $q_uh_u/(6t)$, and $\mu$ is
symmetric in its two arguments.

On $N(u)$ define the transition kernels
\begin{equation}
 P_u(v,z)=\frac{\mathbf 1_{\{vz\in E(F)\}}q_z}{c_{uv}},\qquad
 Q_u=\frac{I+P_u}{2},\qquad
 K_u^{(j)}=P_uQ_u^j\quad(j\ge0).
 \label{eq:neighborhood-kernels}
\end{equation}
All denominators are positive.  The kernel $P_u$ is reversible with
stationary law $\pi_u$, since
\[
 \pi_u(v)P_u(v,z)=\frac{q_vq_z\mathbf 1_{\{vz\in E(F)\}}}{h_u}
\]
is symmetric in $v,z$.  The same holds for $Q_u$ and $K_u^{(j)}$, which
are polynomials in $P_u$.  Moreover, writing
$K_u^{(j)}=Q_u^jP_u$ and using $c_{uv}\ge\ell^{-1}$ gives
\begin{equation}
 K_u^{(j)}(v,z)\le\ell q_z.
 \label{eq:kernel-density-bound}
\end{equation}

The number of vertices in $N(u)$ can be arbitrarily large even when its
total weight is small.  We therefore measure entropy relative to $q$:
for a probability row $p$ on a finite set with positive weights $q$, put
\[
 \mathcal H_q(p)=\sum_{z:p(z)>0}p(z)\log\frac{q_z}{p(z)}.
\]
This quantity may be negative.  Jensen's inequality bounds it by
$\log q(\operatorname{supp}p)$.  Define
\begin{equation}
 H_j=\EE_{(u,v)\sim\mu}
            \mathcal H_q\bigl(K_u^{(j)}(v,\cdot)\bigr),
 \qquad R=\lceil\ell\rceil.
 \label{eq:walk-entropy}
\end{equation}
The density bound gives $H_j\ge-\log\ell$.  We will bound $H_R$ by
$O_r(\log\ell)$, much less than the immediate bound $H_R\le\ell$.
Only $O_r(\log\ell)$ entropy can then be gained over $R$ steps, so
some step has a small increment.  That increment will control the
distance between the label laws at a triangle corner.

\begin{lemma}
\label{lem:walk-entropy-bound}
With the above definitions,
\begin{equation}
 H_0\ge-\log\ell,\qquad H_R\le O_r(\log\ell).
 \label{eq:walk-entropy-bound}
\end{equation}
\end{lemma}

\begin{proof}
Only the upper bound remains.  Fix $v$.  We construct sets that contain
almost all the mass of the walks starting at $v$, uniformly over the
choice of $u\in N(v)$.  Let
\[
 \theta=\ell^{-6}e^{-\ell},\qquad D_0=N(v),\qquad
 D_{j+1}=D_j\cup\{z:q(N(z)\cap D_j)\ge\theta\}
                      \quad(0\le j<R).
\]
The sets $D_j$ lie in the ball of radius $j+1$ about $v$.  Double counting
the edges between $D_j$ and its newly included vertices yields
\[
 \theta q(D_{j+1}\setminus D_j)\le e^\ell q(D_j).
\]
Consequently,
\begin{equation}
 q(D_R)\le e^\ell(1+\ell^6e^{2\ell})^R\le e^{20\ell^2}.
 \label{eq:walk-support-mass}
\end{equation}
For the last inequality, use $6\log\ell\le\ell$ and $R\le2\ell$ for
$\ell\ge32$.

We also need the part condition in \eqref{eq:local-cross}.  For a fixed
part $D'$, every vertex in $(D_{j+1}\setminus D_j)\cap D'$ has at least
$\theta$ weight of neighbors in $D_j$.  Each vertex of $D_j$ has at
most one neighbor in $D'$, so
\[
 \theta\bigl|(D_{j+1}\setminus D_j)\cap D'\bigr|\le q(D_j).
\]
Thus the weighted inclusion threshold controls an unweighted count in
each part.  The initial set $N(v)$ has at most one vertex in $D'$.  It follows from
\eqref{eq:walk-support-mass} that
\begin{equation}
 |D_R\cap D'|\le1+R\theta^{-1}e^{20\ell^2}
       \le e^{40\ell^2}\le P_0.
 \label{eq:walk-support-parts}
\end{equation}
The radius condition also holds, since $R+1\le R_0$.

For every $u\in N(v)$, the leakage outside these sets satisfies
\begin{equation}
 K_u^{(j)}(v,N(u)\setminus D_j)\le j\ell^{-4}
                          \qquad(0\le j\le R).
 \label{eq:walk-leakage}
\end{equation}
At $j=0$, the row is supported on $N(u)\cap N(v)$.  For the induction
step, fix $z\in N(u)\setminus D_{j+1}$.  Since $D_j\subseteq D_{j+1}$,
mass arriving at $z$ from $D_j$ uses an off-diagonal transition of $Q_u$.
By \eqref{eq:kernel-density-bound} and the off-diagonal bound
$Q_u(y,z)\le\ell q_z$, this mass is at most
\[
 \sum_{y\in N(z)\cap D_j}\ell q_y\,\ell q_z
 \le \ell^2q_z\theta
 =q_z\ell^{-4}e^{-\ell}.
\]
Summing over these $z$ gives at most $\ell^{-4}$, since $q(N(u))\le
e^\ell$.  The mass arriving from outside $D_j$ is at most the preceding
leakage, proving \eqref{eq:walk-leakage}.

Although $D_R$ can have large total weight, its intersection with most
neighborhoods relevant to the triangle law is small.  Define
\[
 J_v=\{u\in N(v):q(N(u)\cap D_R)>\ell^8\}.
\]
Apply \eqref{eq:local-cross} to $N(v)$ and $D_R$, with $h=20\ell^2$
and $\eta=2$.  The mass, radius, and part conditions were just verified;
also $h\le20\ell_0^2$.  As $S_v\ge\ell^{-1}$, the exponential error
can be absorbed, giving
\[
 \ell^8q(J_v)\le e_q(N(v),D_R)\le O_r(\ell^2S_v),
 \qquad q(J_v)\le O_r(\ell^{-6}S_v).
\]
This controls $q$-mass, whereas the conditional law given $v$ in
\eqref{eq:triangle-laws} also weights $u$ by $c_{uv}$.  A second use of
\eqref{eq:local-cross}, now on $J_v,N(v)$ with $h=\ell$ and $\eta=2$,
makes precisely this conversion:
\begin{equation}
 \begin{aligned}
 \PP_\mu(u\in J_v\mid v)
 &=\frac{\sum_{u\in J_v}q_uc_{uv}}{h_v}
   =\frac{e_q(J_v,N(v))}{h_v}\\
 &\le O_r\!\left(
       \frac{\ell q(J_v)+e^{-3\ell}}{S_v/\ell}\right)
   \le O_r(\ell^{-4}).
 \end{aligned}
 \label{eq:exceptional-neighborhood-law}
\end{equation}
Here the sets are contained in $N(v)$, so the support conditions hold
with radius one and at most one vertex per part.

For $u\notin J_v$, split the row $K_u^{(R)}(v,\cdot)$ according to
membership in $D_R$.  Its outside probability is at most $R\ell^{-4}$
by \eqref{eq:walk-leakage}; the weights of the two possible supports
are at most $\ell^8$ and $e^\ell$, respectively.  The entropy of this
mixture is therefore at most
\[
 \log2+8\log\ell+(R\ell^{-4})\ell.
\]
Indeed, conditioning on the two supports adds their binary entropy,
which is at most $\log2$, and the conditional entropies are bounded by
the logarithms of their support weights.  Empty supports contribute
nothing.  For $u\in J_v$, use the bound $\ell$.  Averaging first over
$u$ conditional on $v$, using
\eqref{eq:exceptional-neighborhood-law}, and then over $v$ proves
$H_R=O_r(\log\ell)$.
\end{proof}

\subsection{Coupled labels and the split graph}
\label{subsec:coupled-splitting}

\begin{proof}[Proof of Lemma~\ref{lem:triangle-splitting}]
We first turn the entropy bound into similarity of the label laws at a
triangle corner.  For two probability rows $p,a$ on the same finite set,
write
\[
 \KL(p\|a)=\sum_{z:p(z)>0}p(z)\log\frac{p(z)}{a(z)},\qquad
 \TV(p,a)=\frac12\sum_z|p(z)-a(z)|,
\]
with $\KL(p\|a)=+\infty$ when necessary.  We shall use
\begin{equation}
 \TV(p,a)\le\sqrt{\KL(p\|a)}.
 \label{eq:entropy-total-variation}
\end{equation}
For completeness, let $\rho=\sum_z\sqrt{p(z)a(z)}$.  Jensen's
inequality gives $\KL(p\|a)\ge-2\log\rho\ge2(1-\rho)$, while
Cauchy--Schwarz gives $\TV(p,a)^2\le1-\rho^2\le2(1-\rho)$.
The case of infinite divergence is immediate.

Since $K_u^{(j+1)}=Q_uK_u^{(j)}$, its row at $v$ is the mixture of
the rows $K_u^{(j)}(z,\cdot)$ with $z\sim Q_u(v,\cdot)$.  The entropy
of a mixture minus the average entropy of its constituent rows is the
average of their divergences from the mixture.  The $\log q_z$ terms
cancel in this identity.  Averaging over $\mu$, and using stationarity
of $\pi_u$ for $Q_u$ to identify the average constituent entropy as
$H_j$, gives
\begin{equation}
 H_{j+1}-H_j
 =\EE_{\substack{(u,v)\sim\mu\\z\sim Q_u(v,\cdot)}}
   \KL\bigl(K_u^{(j)}(z,\cdot)\,
             \big\|\,K_u^{(j+1)}(v,\cdot)\bigr)\ge0.
 \label{eq:entropy-increment}
\end{equation}
Lemma~\ref{lem:walk-entropy-bound} and telescoping over $R$ steps now
give an index $0\le j<R$ such that
\[
 H_{j+1}-H_j\le O_r((\log\ell)/\ell).
\]
Fix this index and abbreviate $K_u=K_u^{(j)}$.  Using the row
$K_u^{(j+1)}(v,\cdot)$ as an intermediary in the triangle inequality
for total variation yields
\begin{equation}
 \EE_{\substack{(u,v)\sim\mu\\z\sim P_u(v,\cdot)}}
       \TV\bigl(K_u(v,\cdot),K_u(z,\cdot)\bigr)
 \le2\sqrt{H_{j+1}-H_j}
 \le O_r\!\left(\sqrt{\frac{\log\ell}{\ell}}\right).
 \label{eq:corner-row-agreement}
\end{equation}
To check the factor two, the two expected distances to the intermediary
sum to twice the expected distance from the $z$-row under
$Q_u=(I+P_u)/2$.  Apply \eqref{eq:entropy-total-variation}, then
Cauchy--Schwarz, and finally \eqref{eq:entropy-increment}.  The law on
the left of \eqref{eq:corner-row-agreement} is exactly the ordered
triangle law, because
\[
 \mu(u,v)P_u(v,z)=\frac{q_uq_vq_z}{6t}
 \quad\text{when $u,v,z$ form a triangle}.
\]

At each $u$, give all incidences $(u,v)$ labels $Y_{u,v}\in N(u)$,
with marginal law $K_u(v,\cdot)$, by a simultaneous coupling.  We use
the disagreement bound of Angel and Spinka~\cite[Theorem~2]{AngelSpinka}; the
following finite-set construction includes the proof needed here.
For any finite family of laws on a finite set, take one shared sequence
of independent proposals, each consisting of a uniformly chosen label
and an independent uniform mark in $[0,1]$.  Each law accepts the first
proposal whose mark is no greater than its probability of that label.
The accepted label has the prescribed law, and every law accepts almost
surely.  For two laws at total variation distance $d$, the first
proposal accepted by either is accepted by both with probability
$(1-d)/(1+d)$: the sums of the pointwise minimum and maximum of their
probabilities are $1-d$ and $1+d$.  Their labels thus disagree with
probability at most $2d/(1+d)\le2d$.  A single proposal sequence gives
this bound simultaneously for every pair of laws.  Applying the
construction to the rows of $K_u$ and using
\eqref{eq:corner-row-agreement} bounds the average probability of a
disagreement at a triangle corner by
$O_r(\sqrt{(\log\ell)/\ell})$.  Couplings at distinct vertices may
be chosen independently.

Agreement controls triangle retention.  To control neighborhood
weights, call an incidence $(u,v)$ with label $y=Y_{u,v}$
\emph{qualified} if
\begin{equation}
 K_u(y,v)\ge q_ve^{-\sqrt\ell}.
 \label{eq:reverse-qualification}
\end{equation}
This uses the reverse transition: at a copy indexed by $y$, summing
$K_u(y,v)$ over its possible neighbors will sum a single probability
row.  It remains to show that this qualification rarely discards an
incidence.

Choose $(u,v)\sim\mu$ and then $y\sim K_u(v,\cdot)$.  Conditional
on $u$, the joint law of $(v,y)$ is $\pi_u(v)K_u(v,y)$, which is
symmetric by reversibility.  In particular, the reverse entry is
positive almost surely, and exchanging $v$ and $y$ gives
\begin{equation}
 \EE\log\frac{q_v}{K_u(y,v)}
  =\EE\log\frac{q_y}{K_u(v,y)}=H_j\le H_R.
 \label{eq:reverse-entropy}
\end{equation}
The random logarithm is at least $-\log\ell$ by
\eqref{eq:kernel-density-bound}.  Add $\log\ell$ and apply Markov's
inequality to this nonnegative variable.  Lemma~\ref{lem:walk-entropy-bound}
then yields
\begin{equation}
 \PP\bigl((u,v)\text{ is not qualified}\bigr)
 \le\frac{H_j+\log\ell}{\sqrt\ell+\log\ell}
 \le O_r\!\left(\frac{\log\ell}{\sqrt\ell}\right).
 \label{eq:qualification-failure}
\end{equation}

For each $u$ make copies $(u,y)$ indexed by $y\in N(u)$, all with
weight $q_u$.  Retain $uv$ if its two incidences qualify, and join
$(u,Y_{u,v})$ to $(v,Y_{v,u})$.  This defines a simple graph, since
each old edge contributes at most one edge and its endpoints project
to distinct vertices.  The neighbors of $(u,y)$ have total weight at
most
\begin{equation}
 \sum_{\substack{v\in N(u)\\
                 Y_{u,v}=y,\ uv\text{ retained}}}q_v
 \le e^{\sqrt\ell}\sum_{v\in N(u)}K_u(y,v)
 =e^{\sqrt\ell},
 \label{eq:split-neighborhood-bound}
\end{equation}
by \eqref{eq:reverse-qualification}.  Figure~\ref{fig:vertex-splitting}
illustrates how agreement or disagreement of the corner labels
determines whether a triangle survives.

\begin{figure}[t]
\centering
\begin{tikzpicture}[x=1.1cm,y=1.1cm,
 vertex/.style={circle,fill=black,inner sep=1.6pt},
 blueedge/.style={draw=blue!65!black,line width=0.9pt},
 orangeedge/.style={draw=orange!85!black,line width=0.9pt},
 every node/.style={font=\small}]
 \node[vertex,label=left:$u$] (u) at (0,0) {};
 \node[vertex,label=above:$a$] (a) at (1.5,1.35) {};
 \node[vertex,label=right:$b$] (b) at (2.9,0.75) {};
 \node[vertex,label=below:$c$] (c) at (1.5,-1.35) {};
 \node[vertex,label=right:$d$] (d) at (2.9,-0.75) {};
 \draw[blueedge] (u)--node[above left]{$y$}(a)--(b);
 \draw[blueedge] (u)--node[pos=.32,above]{$y$}(b);
 \draw[orangeedge] (u)--node[below left]{$z$}(c)--(d);
 \draw[orangeedge] (u)--node[pos=.32,below]{$z$}(d);
 \draw[gray,line width=0.6pt] (a)--(c);
 \node at (1.45,-2.05) {before splitting at $u$};
 \draw[-{Stealth[length=2mm]},thick] (3.5,0)--(4.3,0);
 \node[vertex,label=left:{$(u,y)$}] (uy) at (5.2,0.85) {};
 \node[vertex,label=left:{$(u,z)$}] (uz) at (5.2,-0.85) {};
 \node[vertex,label=above:$a$] (aa) at (6.7,1.35) {};
 \node[vertex,label=right:$b$] (bb) at (8.1,0.75) {};
 \node[vertex,label=below:$c$] (cc) at (6.7,-1.35) {};
 \node[vertex,label=right:$d$] (dd) at (8.1,-0.75) {};
 \draw[blueedge] (uy)--(aa)--(bb)--(uy);
 \draw[orangeedge] (uz)--(cc)--(dd)--(uz);
 \draw[gray,line width=0.6pt] (aa)--(cc);
 \node at (6.5,-2.05) {each copy retains weight $q_u$};
\end{tikzpicture}
\caption{A local view of the splitting construction. All displayed
incidences are assumed to qualify, and labels at vertices other than $u$
are suppressed. These are auxiliary incidence labels, not colors in the
correspondence assignment. The incidences $ua,ub$ have label $y$, while $uc,ud$ have
label $z\ne y$. The triangles $uab$ and $ucd$ remain triangles after
splitting at $u$; $uac$ does not. Each old edge is used once, and each
copy of $u$ keeps its original weight. The coupling makes unequal labels
at a triangle corner rare in weighted average.}
\label{fig:vertex-splitting}
\end{figure}


An old triangle survives with its original mass if all six incidences
qualify and the two labels agree at each of its three corners.  In the
ordered triangle law, each incidence has marginal $\mu$, and each
corner has the law used in \eqref{eq:corner-row-agreement}.  A union
bound, followed by averaging with the triangle weights, shows that the
expected fraction of triangle mass lost is at most
\[
 O_r\!\left(\frac{\log\ell}{\sqrt\ell}
              +\sqrt{\frac{\log\ell}{\ell}}\right)
       =O_r\!\left(\frac{\log\ell}{\sqrt\ell}\right).
\]
Choose a realization retaining at least the expected lower bound.  The
projection to $F$ is a homomorphism injective on directed edges, and
composing it with the old map to $H_b$ preserves that property.  All
vertex weights and inherited data remain unchanged.  This proves
\eqref{eq:split-retention} and the lemma.
\end{proof}

\subsection{Iteration and completion of the adjustment lemma}
\label{subsec:triangle-iteration}

\begin{proof}[Proof of Proposition~\ref{prop:triangle-bound}]
Apply the deletions of Subsection~\ref{subsec:triangle-preparation} to
all levels, then Lemma~\ref{lem:triangle-splitting} to each nonempty
level graph.  Empty level graphs need no construction.  The resulting
graphs have neighborhood weights at most $e^{\sqrt\ell}$ and still
have homomorphisms into their original $H_b$ injective on directed
edges.  Combining the additive losses from deletion with
\eqref{eq:split-retention} gives
\begin{equation}
 T_{\mathrm{new}}\ge
 \left(1-C_2\frac{\log\ell}{\sqrt\ell}\right)T
       -C_3\frac{\log\ell}{\ell}M_x,
 \label{eq:triangle-iteration}
\end{equation}
where the constants depend only on $r$.  If the coefficient of $T$ is
negative, the inequality follows simply from $T_{\mathrm{new}}\ge0$.

Choose a sufficiently large constant $X\ge32$, depending only on $r$.
Starting at $\ell_0$, perform the reduction while $\ell>X$, replacing
$\ell$ by $\sqrt\ell$ each time.  Both losses in
\eqref{eq:triangle-iteration} are summable uniformly in the initial
value $\ell_0$.  To verify this explicitly, let
$p(s)=(\log s)/\sqrt s$.  This function is decreasing for sufficiently
large $s$, and
\[
 \frac{p(s^2)}{p(s)}=\frac2{\sqrt s}.
\]
In reverse order, the parameters at which reductions are performed
start above $X$ and are successively squared.  Increasing $X$ therefore
makes the sum of $C_2p(\ell)$ at most $1/2$.  The sum of
$C_3(\log\ell)/\ell$ is bounded by a constant depending only on $r$.
Iterating \eqref{eq:triangle-iteration}, and using
$\prod_j(1-a_j)\ge1-\sum_ja_j$ for these nonnegative losses, gives
\begin{equation}
 T_{\mathrm{final}}\ge\frac12T_0-O_r(M_x).
 \label{eq:iterated-triangle-retention}
\end{equation}

At the final stage every neighborhood has weight at most $e^X$.
Every triangle still has a middle vertex.  At a middle vertex $a$,
the incident triangle mass, including the factor $b$, is at most
$bq_aq(N(a))^2/2$.  Hence \eqref{eq:middle-edge-charge} gives
\[
 T_{\mathrm{final}}
 \le\frac{e^X}{2}
      \sum_b b\sum_{\substack{a\in V(F_b)\\a\text{ middle}}}
                     q_aq(N(a))
 \le e^XM_x.
\]
This estimate also applies if $\ell_0\le X$ and no reduction is needed.
Together with \eqref{eq:iterated-triangle-retention}, it proves
$T_0=O_r(M_x)$.
\end{proof}

\begin{proof}[Completion of the proof of Lemma~\ref{lem:adjustment}]
Proposition~\ref{prop:triangle-bound} bounds the contribution from
triangles all of whose vertices are good by $O_r(M_x)$.  The contribution
from triangles meeting a bad vertex is, by \eqref{eq:bad-triangles}, at
most $O(\Lambda^2)\sum_iw_iy_i$.  Taking $B=B(r)$ sufficiently large
and then increasing the lower threshold on $\Lambda$ gives
\[
 \sum_i x_it_i\le B\sum_i x_id_i+\Lambda^4\sum_iw_iy_i,
\]
which is the dual inequality \eqref{eq:dual-target}.  The separation
argument in Section~\ref{sec:adjustment} now supplies the finite family
of operations and rates in \eqref{eq:adjustment}.  The constants are
chosen in order: first the multiplier bound $A$ and the local estimate
constant $C_1$, then the reduction constants and $X$, and finally $B$
and the threshold on $\Lambda$.  Every choice depends only on $r$, as
required.
\end{proof}
